\chapter{Conclusion and Future Work}

\section{Conclusion}
This thesis introduced and evaluated AI-ActivEdu, a cloud-based platform for H5P interactive video authoring designed around educator workflows. The system integrates timeline-based authoring, AI-assisted drafting, standards-compatible packaging, and operational controls suitable for institutional deployment.

Comparative evaluation demonstrated substantial usability and productivity improvements over plugin-centric alternatives. These outcomes indicate that architecture and UX decisions can be as critical as algorithmic capability in educational technology adoption.

\section{Contributions Recap}
The principal contributions are:
\begin{enumerate}[leftmargin=2em]
\item A reproducible, service-oriented architecture for interactive video authoring.
\item A constrained AI pipeline with schema validation and human review.
\item Evidence of measurable workflow gains in task completion, time, and satisfaction.
\item A practical security baseline aligned with OWASP-oriented controls.
\end{enumerate}

\section{Future Work}
Future development priorities include:
\begin{itemize}[leftmargin=2em]
\item Client-side video editing features (trim, split, crop, mute) with FFmpeg WASM.
\item Broader question-type coverage and rubric-aligned analytics.
\item Multi-language authoring support with translation-aware review.
\item Deeper LMS telemetry and adaptive feedback recommendations.
\item Larger multi-institution evaluations for external validity.
\end{itemize}

\section{Closing Statement}
The project demonstrates that educator-centered design, modern web architecture, and responsibly integrated AI can significantly lower barriers to interactive content creation in higher education.
